\documentclass[11pt]{article}
\usepackage[margin=1in]{geometry}
\usepackage{amsmath,amssymb,amsthm}
\usepackage{mathrsfs}
\usepackage[hidelinks]{hyperref}
\usepackage{enumitem}
\setlength{\emergencystretch}{3em}

\newtheorem{theorem}{Theorem}
\newtheorem{lemma}[theorem]{Lemma}
\newtheorem{proposition}[theorem]{Proposition}
\newtheorem{corollary}[theorem]{Corollary}
\newtheorem{conjecture}[theorem]{Conjecture}
\newtheorem*{lemmaA}{Lemma A}
\theoremstyle{definition}
\newtheorem{definition}[theorem]{Definition}
\theoremstyle{remark}
\newtheorem{remark}[theorem]{Remark}

\newcommand{\R}{\mathbb{R}}
\newcommand{\Z}{\mathbb{Z}}
\newcommand{\N}{\mathbb{N}}
\newcommand{\Q}{\mathbb{Q}}
\newcommand{\ee}{\mathrm{e}}
\newcommand{\dist}{\operatorname{dist}}
\newcommand{\dimH}{\dim_{\mathrm{H}}}
\newcommand{\Lam}{\Lambda}

\title{Complete spectrum discretization, exact Hausdorff gauges,\\
and dimension phase transitions for Diophantine approximation\\
along the convergent denominators of \texorpdfstring{$e$}{e}}
% TODO before submission: replace with the real author name(s) and affiliation(s).
\author{[Author name]\\ \relax[Affiliation]}
\date{Integrated draft, November 23, 2026}

\begin{document}
\maketitle

\begin{abstract}
Let $(Q_j)_{j\ge1}$ be the selected convergent denominators of $\ee$
immediately preceding the partial quotients $2j$
($Q_1=1,Q_2=7$, $Q_{j+1}=(4j+2)Q_j+Q_{j-1}$), let $\|y\|=\dist(y,\Z)$, and
define the weighted approximation exponent
$\Lam(x)=\limsup_{j\to\infty} j\,\|Q_j x\|$. We prove that the full spectrum is
completely discrete:
$\Lam(\R)=\frac14\Z_{\ge0}\cup\{\infty\}=\{0,1/4,1/2,3/4,1,\dots\}\cup\{\infty\}$,
so every open interval $(k/4,(k+1)/4)$ ($k\in\Z_{\ge0}$) is an unconditional
empty spectral gap, and $\Lam(x)=0$ exactly on the countable group
$L=\Z+\Z\frac{\ee+1}{2}$.

Our main metric results measure the critical levels and their fine structure.
For an integer $k\ge1$ let $h_k$ be the gauge with $h_k(1/Q_n)=(2k+1)^{-n}$,
interpolated linearly. On the circle $\R/\Z$, the exact level $\{\Lam=k/4\}$
has infinite but $\sigma$-finite $\mathcal H^{h_k}$ measure, while
$\{\Lam<k/4\}$ is $\mathcal H^{h_k}$-null. The gauge $h_k$ is exact: the level
has infinite $\mathcal H^{h_k}$ measure in every nonempty open arc, every gauge
of smaller order gives $\{\Lam\le k/4\}$ measure zero, and for every gauge of
larger order $\{\Lam=k/4\}$ is not $\sigma$-finite. For arbitrary gauges,
including oscillating ones, the three measure regimes are determined by
$\liminf_n (2k+1)^n g(1/Q_n)$, or by $\liminf_n 2^n g(1/Q_n)$ for the
unsigned exact-limit set. For finite-state tail constraints the base is
the largest Perron root among components supporting infinitely many
extreme defects. Furthermore, at each
positive level $c=k/4$ ($k\ge1$) we establish a sharp three-way splitting
between the limsup level (digit alphabet of size $2k+1$, critical gauge base
$2k+1$), the unsigned exact limit set $\{x:j\|Q_jx\|\to k/4\}$ (universal binary
alphabet $\{\pm k\}$, critical gauge base $2$ \emph{independently} of
$k\ge1$), and each signed exact limit set
$\{x:j(Q_jx-\mathrm{round}(Q_jx))\to\pm k/4\}$ (which collapses to a single
countable coset of $L$). More generally, for any polynomial-growth recurrence
$Q_{j+1}=(cj^r+d)Q_j+Q_{j-1}$ ($r,c\in\Z_{\ge1}$, $d\in\Z$), the critical
exponent $\Lam^{(r,c,d)}(x)=\limsup_{j\to\infty}j^r\|Q_jx\|$ has full
spectrum $\frac1c\Z_{\ge0}\cup\{\infty\}$, intrinsic rank-$2$ zero-level
lattice $L^{(r,c,d)}$, critical gauges $h_k^{(r,c,d)}(1/Q_n)=(2k+1)^{-n}$
with $1/r$-scaled logarithmic order and positive-level three-way splitting, and
subcritical Hausdorff dimension phase transition $1-\alpha/r$ ($0\le\alpha<r$).
Even more generally, for \emph{any} growing-quotient recurrence
$Q_{j+1}=A_jQ_j+Q_{j-1}$ ($A_j\in\Z_{\ge2}$, $A_j\to\infty$), the normalized
approximation functional $\Lam_G(x)=\limsup_{j\to\infty}A_j\|Q_jx\|$ agrees
identically with the Kraaikamp--Liardet error ratio, has full discrete spectrum
$\Z_{\ge0}\cup\{\infty\}$ with zero-level lattice $L_G=\frac{1}{g}(\Z+\Z\theta)$
($g=\gcd(Q_{J_0-1},Q_{J_0})$), universal critical and binary gauges on $\R$ and
$\R/\Z$, subcritical Hausdorff dimension $(1-\tau)/(1+\tau\gamma_{\sup})$ for
$\Lam_{G,\tau}(x)=\limsup_{j\to\infty}A_j^\tau\|Q_jx\|$ ($0\le\tau<1$, where
$\gamma_{\sup}=\limsup_{j\to\infty}\frac{\log A_j}{\log Q_j}\in[0,\infty]$),
and a supercritical ($\tau>1$) secondary phase transition at
$\tau_{\mathrm{crit}}=1+1/\gamma_{\sup}$ where finite sublevel sets jump from
the rank-$2$ lattice $L_G$ to the strict rank-$1$ rational sublattice
$\frac{1}{g}\Z\subsetneq L_G$.

Next, every finite level set $\{\Lam\le c\}$ has Hausdorff dimension $0$.
Hence $\{\Lam<\infty\}$ has dimension $0$, and $\Lam=\infty$ Lebesgue-almost
everywhere (with Hausdorff dimension $1$). For $0\le\alpha<1$, the sublevel set $\{\Lam_\alpha\le c\}$ for $c\ge0$
has Hausdorff dimension $1-\alpha$. The same holds for the set where
$j^\alpha\|Q_jx\|\to c$ for $c\ge0$ if $\alpha>0$, and for
$c\in[0,1/2]$ if $\alpha=0$. The unweighted case also follows from work
of Bugeaud, Liao and Rams. For ordinary continued fractions,
we determine the Hausdorff dimensions of Hermite window-count constraints,
show that eventual periodicity of the partial quotients modulo $m$ ($m\ge2$) has
the same Hausdorff dimension as the fixed residue alphabet $F_{m,1}$ (strictly
between $357/512$ and $179/256$ for $m=4$), prove the global strict separation
$\dimH F_{m,m}<\dimH F_{m,1}=\dimH E_m$ for \emph{all} moduli $m\ge2$ via an
exact resolvent identity, and establish the two-term asymptotic expansion
$\dimH F_{m,r}=\frac12+\frac{C_r}{2m}-\frac{C_r^2\log m}{2m^2}+O(m^{-2})$ as
$m\to\infty$ for fixed $r\ge1$ (where $C_r=\sum_{k=1}^\infty 1/q_k(r)$ is the
reciprocal metallic-continuant series and $C_1=\psi_{\mathrm{Fib}}=\sum_{k=1}^\infty 1/F_k\approx3.3598856662$ is the reciprocal Fibonacci constant), while
$\dimH F_{m,m}=\frac12+\frac{1}{2m}-\frac{\log m}{2m^2}+O(m^{-2})$. More generally,
for finite multi-residue alphabets $F_{m,S,B}$ ($S\subset\Z_{\ge1}$,
$B\subset\Z_{\ge0}$, $d=|S|+|B|\ge1$), Theorem~\ref{thm:multi-residue-asymptotic}
gives a self-contained, Bowen-free proof of the classical identity
$\dimH E_S=s_0(S)$ and establishes the universal phase transition
$\dimH F_{m,S,B}\to\max(1/2,s_0(S))$, a three-term
Euler--Mascheroni subcritical expansion to $O((\log m)^2/m^3)$ with coefficient
$\frac12 d^2 C_S(\gamma C_S-C_S')$, and the supercritical power-law convergence
rate $\Theta(m^{-2s_0(S)})$ for $s_0(S)>1/2$.

The classification of $\{\Lam=0\}$, the initial and full spectral-gap ladder,
Theorem~\ref{thm:gauge},
Corollaries~\ref{cor:exact-gauge} and~\ref{cor:liminf-gauges},
Remark~\ref{rem:gauge-scale}, all three parts (i)--(iii) of
Theorem~\ref{thm:trichotomy} on both $\R$ and $\R/\Z$ (\texttt{UnitAddCircle},
including the binary exact-limit gauge regime and the signed-coset collapse
$E_{\mathrm{lim}}^\pm(k/4)=x_k^\pm+L$), the dimension-zero theorem for all
finite sublevels (and full Lebesgue measure of $\{\Lam=\infty\}$), both the upper
and lower bounds for the subcritical dimension $1-\alpha$, the complete
general polynomial-growth recurrence theory ($r\ge1$)---Lemmas~\ref{lem:poly-growth}--\ref{lem:poly-lattice},
Theorems~\ref{thm:general-rec}--\ref{thm:poly-phase-transition},
Corollary~\ref{cor:poly-exact-gauge}, and Remark~\ref{rem:poly-gauge-scale}---and
its universal extension to arbitrary growing-quotient recurrences
$Q_{j+1}=A_jQ_j+Q_{j-1}$ (Theorems~\ref{thm:growing-rec-spectrum}
and~\ref{thm:growing-rec-dim-phase}, and
Corollary~\ref{cor:growing-rec-three-regimes}), together
with all Hausdorff-dimension statements of Theorem~\ref{thm:windows} (both the
exact $\dimH=1$ prescribed-parity statement and the exact $\dimH=1/2$ upper and lower bounds for eventually bounded and prescribed
increments), together with all of Theorems~\ref{thm:short-ae} and~\ref{thm:no-short-dim}
(both the almost-everywhere logarithmic short-window frequency, the
Hausdorff-dimension identity $\dimH=1$, and the Lebesgue-nullity deduction
pointwise for mathlib's regular continued fractions), all parts~(1)--(3) of
Theorem~\ref{thm:finite-state-tail} unconditionally (under distinct outgoing edge
labels at each vertex, eliminating the Perron--Frobenius eigenvector hypothesis via
a Brouwer-free Collatz--Wielandt minimax proof in \texttt{PerronFrobenius.lean}, $13$
declarations audited in \texttt{AuditPerronFrobenius.lean}), as well as the
general-modulus ($m\ge2$) uniform $24$-Lipschitz coordinate-shift sandwich,
eventual-periodicity dimension equality, unconditional transfer-operator
Hausdorff-dimension criterion (bypassing Bowen-pressure axioms), strict lower bound
$\dimH F_{m,r}>1/2$, upper bound $\dimH F_{4,1}\le 4/5$, step-function certificate
soundness for Theorems~\ref{thm:residues} and~\ref{thm:residues-general}, and the
complete resolvent identity, global strict separation
$\dimH F_{m,m}<\dimH F_{m,1}=\dimH E_m$ for all $m\ge2$, and two-term
reciprocal-Fibonacci/metallic asymptotic expansions of Theorem~\ref{thm:residues-asymptotic},
as well as the multi-residue phase transition, Euler--Mascheroni three-term
asymptotics, and supercritical $\Theta(m^{-2s_0(S)})$ rate of
Theorem~\ref{thm:multi-residue-asymptotic} (\texttt{ResidueMultiAsymptotic.lean}, $24$
declarations audited in \texttt{AuditResidueMultiAsymptotic.lean}), the closed
brackets of Table~\ref{tab:residues-general} (\texttt{ResidueCertTable.lean}, $27$
declarations audited in \texttt{AuditResidueCert.lean}), and the strict open-interval
inequalities of all $20$ bounds in Table~\ref{tab:residues-general}
(\texttt{ResidueCertStrictOpen.lean}, $22$ declarations audited in
\texttt{AuditResidueCertStrictOpen.lean})---are machine-checked in Lean~4 using only
the base axioms \texttt{propext}, \texttt{Classical.choice} and \texttt{Quot.sound}
(Section~\ref{sec:verification}: $937$ main audited declarations across $134$ Lean
files in \texttt{Audit.lean} plus $156$ additional declarations across eight
standalone audit modules spanning $356$ total Lean files, including the companion
modules \texttt{MuEEqTwo.lean} [$9$], \texttt{GeneralAlphabetFloor.lean} [$15$],
\texttt{SupercriticalConstant.lean} [$6$ conditional declarations], and
\texttt{FullConvergentSpectrum.lean} [$40$] of spin-off manuscripts~2D and~2E).
None of Theorem~\ref{thm:finite-state-tail} or Table~\ref{tab:residues-general}
requires external assumptions; only the analytic simple-pole residue verification in
Theorem~\ref{thm:multi-residue-asymptotic}'s supercritical constant / critical
boundary (treated in spin-off Paper~2D), the general even-period signed Hurwitz
recurrence (spin-off Paper~2E), and the irrationality of $e+\pi$ remain outside the
Lean kernel.
The classification and the first gap are a subsequence analogue of a criterion
of Kraaikamp and Liardet. The work is a spin-off of a project on the
irrationality of $e+\pi$ and does \emph{not} resolve that open problem.
\end{abstract}

\section{Introduction}

This article studies the metric theory of Diophantine approximation along a
\emph{sparse, super-exponentially growing} subsequence of denominators: the
selected convergent denominators $(Q_j)_{j\ge1}$ of $\ee$ (notation in the
Setup below). We first show in Theorem~\ref{thm:full-spectrum} that the integer
three-term recurrence for $(Q_j)$ forces complete discretization of the weighted
approximation spectrum:
\[
\Lam(\R)=\frac14\Z_{\ge0}\cup\{\infty\}
=\left\{0,\frac14,\frac12,\frac34,1,\dots\right\}\cup\{\infty\},
\]
creating an infinite family of equidistant empty spectral gaps
$(k/4,(k+1)/4)$ of width $1/4$ and forcing the sublevel sets
$E_c=\{x:\Lam(x)\le c\}$ to be constant across each gap
($E_c=E_{k/4}$ for $c\in[k/4,(k+1)/4)$, Corollary~\ref{thm:sub}).

Our main metric result at the critical weight, Theorem~\ref{thm:gauge},
measures each positive quantized level of $\Lam$: for each integer $k\ge1$
there is an explicit Hausdorff gauge $h_k$ with $h_k(1/Q_n)=(2k+1)^{-n}$, built
from the recurrence itself, for which the exact level $\{\Lam=k/4\}$ has
infinite but $\sigma$-finite Hausdorff measure $\mathcal H^{h_k}$ on the circle
$\R/\Z$, while the strict sublevel $\{\Lam<k/4\}$ is $\mathcal H^{h_k}$-null.
This differs both from the $h$-sets of Cabrelli, Mendivil, Molter and
Shonkwiler \cite[\S5, Theorem~3]{CMMS}, whose gauge gives a finite positive
measure, and from the classical Lagrange level at value $3$ studied by Moreira,
Erazo and Angelini \cite{MEA}, which has either zero or non-$\sigma$-finite
measure for every dimension function. Corollary~\ref{cor:exact-gauge} shows
that $h_k$ is the exact critical gauge of the level $k/4$: smaller gauges give
measure zero even on the non-strict sublevel $\{\Lam\le k/4\}$, larger gauges
destroy $\sigma$-finiteness on $\{\Lam=k/4\}$, and the $\mathcal H^{h_k}$
measure of $\{\Lam=k/4\}$ is infinite in every nonempty open arc.
Corollary~\ref{cor:liminf-gauges} classifies every nondecreasing gauge by
its lower limit along the recurrence scales, even when the normalized
cost oscillates. Theorem~\ref{thm:finite-state-tail} shows how this
classification changes under constraints on consecutive defects.
The largest Perron root is taken only over components in which a path
can revisit an extreme defect infinitely often.
Furthermore, Theorem~\ref{thm:trichotomy} establishes a sharp three-way
splitting at each positive spectral level $k/4$ ($k\ge1$) comparing the limsup
level $\{\Lam=k/4\}$ (gauge base $2k+1$), the unsigned exact limit set
$\{x:j\|Q_jx\|\to k/4\}$ (universal critical gauge base $2$ for \emph{all}
$k\ge1$), and the signed exact limit sets
$\{x:j(Q_jx-\mathrm{round}(Q_jx))\to\pm k/4\}$ (each of which collapses to a
single countable $L$-coset). Theorems~\ref{thm:general-rec}--\ref{thm:poly-phase-transition}
extend the quantization, reference gauges, three-way splitting and
dimension results to general polynomial-growth recurrences
$Q_{j+1}=(cj^r+d)Q_j+Q_{j-1}$ ($r,c\in\Z_{\ge1}$, $d\in\Z$), replacing the
$\ee$-convergent lattice $L$ by an intrinsic rank-$2$ lattice $L^{(r,c,d)}$
spanned by fundamental solutions and establishing the discretized spectrum
$\frac1c\Z_{\ge0}\cup\{\infty\}$, exact gauges with $1/r$-scaled logarithmic
asymptotics, positive-level trichotomy, and subcritical dimension phase
transition $1-\alpha/r$ ($0\le\alpha<r$). In Section~\ref{sec:growing-rec}
(Theorems~\ref{thm:growing-rec-spectrum} and~\ref{thm:growing-rec-dim-phase},
and Corollary~\ref{cor:growing-rec-three-regimes}), we remove the polynomial
growth hypothesis altogether and establish the complete spectral, gauge, and
Hausdorff-dimension phase diagram for arbitrary integer recurrences
$Q_{j+1}=A_jQ_j+Q_{j-1}$ with $A_j\ge2$ and $A_j\to\infty$, uncovering both the
universal subcritical distortion law $\dimH=(1-\tau)/(1+\tau\gamma_{\sup})$ and
a secondary supercritical phase transition at $\tau_{\mathrm{crit}}=1+1/\gamma_{\sup}$
where the level sets collapse from the rank-$2$ lattice $L_G$ to the rank-$1$
rational sublattice $\frac{1}{g}\Z$.

On the coarse scale of Hausdorff dimension, Theorem~\ref{thm:finite} shows that
every finite level set $E_c$ has Hausdorff dimension $0$, so $\{\Lam<\infty\}$
has dimension $0$ while its complement $\{\Lam=\infty\}$ has full Lebesgue
measure and Hausdorff dimension $1$. At subcritical weights
$\Lam_\alpha(x)=\limsup_{j\to\infty}j^\alpha\|Q_jx\|$ ($0\le\alpha<1$),
Theorems~\ref{thm:upper} and~\ref{conj:weight} give exact Hausdorff dimension $1-\alpha$ for
$\{x:\Lam_\alpha(x)\le c\}$ for every $c\ge0$, and for
$\{x\in[0,1]:j^\alpha\|Q_jx\|\to c\}$ for every admissible target $c$.
For $\alpha=0$ such targets are $0\le c\le 1/2$.
The proof uses a covering bound and a shifted-grid Moran construction;
the case $\alpha=0$ also follows from Bugeaud, Liao and Rams \cite{BLR}, and
both bounds for $0\le\alpha<1$ are machine-checked in Lean~4. Finally, in
Section~\ref{sec:windows-residues}, Theorems~\ref{thm:windows}--\ref{thm:no-short-dim}
determine the Hausdorff dimensions and almost-sure frequency of Hermite
continued-fraction window constraints, Theorem~\ref{thm:residues} proves
that eventual periodicity of ordinary continued-fraction digits modulo $4$ has
the same Hausdorff dimension as the fixed alphabet $1\bmod 4$, bounded strictly
between $357/512$ and $179/256$, Theorem~\ref{thm:residues-general} extends this
to every modulus $m\ge2$ (via a uniform $24$-Lipschitz sandwich, an unconditional
transfer-operator Hausdorff-dimension framework, and exact integer certificates
for $m\in\{2,3,4,5,6\}$), Theorem~\ref{thm:residues-asymptotic} proves the
global strict separation $\dimH F_{m,m}<\dimH F_{m,1}=\dimH E_m$ for all $m\ge2$
together with the exact two-term large-$m$ asymptotic expansion governed by the
reciprocal Fibonacci constant $\psi_{\mathrm{Fib}}=\sum_{k=1}^\infty 1/F_k$ and
its metallic-continuant generalizations $C_r=\sum_{k=1}^\infty 1/q_k(r)$, and
Theorem~\ref{thm:multi-residue-asymptotic} establishes the multi-residue
phase transition $\lim_{m\to\infty}\dimH F_{m,S,B}=\max(1/2,s_0(S))$, a
self-contained Bowen-free proof of the classical Poincar\'e-abscissa dimension
formula for finite alphabets $E_S$,
three-term Euler--Mascheroni subcritical asymptotics to $O((\log m)^2/m^3)$,
and the two-sided supercritical power-law rate $\Theta(m^{-2s_0(S)})$.

These results build on the zero-level classification and initial spectral gap
(Section~\ref{sec:class}), which arose as machinery in a separate formalization
project on the (open) irrationality of $e+\pi$ \cite{Phase1}.

\subsection{Setup}

The continued fraction of $\ee$ is $[2;1,2,1,1,4,1,1,6,\dots]$. Let $P_j/Q_j$
be the convergent immediately preceding the partial quotient $2j$; the first
few are $3/1,19/7,193/71,2721/1001$. The denominators satisfy
\begin{equation}\label{eq:rec}
Q_{j+1}=(4j+2)Q_j+Q_{j-1}\quad(j\ge2),\qquad Q_1=1,\ Q_2=7,
\end{equation}
so $Q_{j+1}/Q_j\sim 4j$ and $\log Q_j\sim j\log j$ (super-exponential growth).
Write $\|y\|=\dist(y,\Z)$ and define
\[
\Lam(x):=\limsup_{j\to\infty} j\,\|Q_j x\|\in[0,\infty].
\]
More generally, for $\alpha\ge0$ put
$\Lam_\alpha(x):=\limsup_{j\to\infty} j^\alpha\|Q_j x\|\in[0,\infty]$,
so that $\Lam=\Lam_1$.
Let $L:=\Z+\Z\frac{\ee+1}{2}$, a countable dense subgroup of $\R$.

\subsection{Relation to prior work}\label{sec:prior}

The \emph{phenomenon} (a spectral gap plus zero-dimensional sub-threshold level
sets) is classical: in the classical Lagrange spectrum, the spectrum below
$\sqrt5$ (and below $3$) is countable (Hausdorff dimension $0$), and Hall's ray
gives a half-line. Moreira \cite{Moreira} studies the Hausdorff dimension of
portions of the classical Markov and Lagrange spectra below a threshold. Lima
and Moreira \cite{LimaMoreira} prove phase transitions for the \emph{dynamical}
Markov and Lagrange spectra of generic horseshoes; the corresponding question
for the classical spectra is left open there. Note that our level sets $E_c$ are
subsets of the space of inputs $x\in\R$.

The closest prior continued-fraction result is a theorem of Kraaikamp and
Liardet \cite{KL}, quoted as Lemma~4.1 in \cite{BDS}. Let $\alpha$ be irrational
with convergents $p_n/q_n$. Then a real $\beta$ lies outside
$\Z+\Z\alpha$ if and only if
$\|q_n\beta\|>\tfrac14\,q_n|q_n\alpha-p_n|$ for infinitely many $n$. This
is a classification together with a gap, with constant $1/4$, for the
normalized quantity $\|q_n\beta\|/(q_n|q_n\alpha-p_n|)$ along \emph{all}
convergent denominators. Theorems~\ref{thm:class} and~\ref{thm:gap} can be
compared with this criterion in two complementary ways:
\begin{itemize}[leftmargin=1.8em]
\item \emph{As a subsequence along the convergents of $\ee$:} Passing from all
convergent denominators of $\ee$ to the subsequence $(Q_j)$ of every third
convergent denominator changes two things. First, the zero group grows from
$\Z+\Z\ee$ to $L$: for $x=(\ee+1)/2$ the quantity $j\|Q_jx\|$ tends to $0$, so
the infinitely many $n$ supplied by Kraaikamp and Liardet for $\alpha=\ee$ lie
outside the subsequence $(Q_j)$. Second, the normalization changes. Since
$2j+1<1/(Q_j|Q_j\ee-P_j|)<2j+2$, one has
$\limsup_{j\to\infty}\|Q_jx\|/(Q_j|Q_j\ee-P_j|)=2\Lam(x)$. Thus the gap
$\Lam\ge1/4$ reads $\limsup\ge1/2$ in their normalization, and the value
$\Lam(x_0)=1/4$ is attained by the explicit point $x_0$ in
Remark~\ref{rem:unweighted}. The two constants $1/4$ therefore refer to
different normalizations.
\item \emph{As the full convergent sequence of $\beta=(\ee+1)/2$:} Put
$R_j=(P_j+Q_j)/2\in\Z$. The fractions $R_j/Q_j$ are the convergents of
$\beta=(\ee+1)/2=[1;1,6,10,14,\dots]$, whose initial convergent pairs are
$(1,1)$ and $(2,1)$ and whose remaining partial quotients are $4j+2$ for
$j\ge1$. Applying the Kraaikamp--Liardet criterion \cite{KL} (in the form
recorded on p.~106 of \cite{BDS}) directly to $\beta$ gives
$\|Q_jx\|>\tfrac14 Q_j|Q_j\beta-R_j|$ infinitely often for $x\notin L$. Since
the continued-fraction error formula gives $jQ_j|Q_j\beta-R_j|\to1/4$, this
black-box application yields $\Lam(x)\ge1/16$ outside $L$, which already
recovers the zero-level classification $\{\Lam=0\}=L$, whereas the direct
three-term recurrence analysis in Theorems~\ref{thm:gap}
and~\ref{thm:full-spectrum} yields the sharp initial gap $1/4$ and the complete
quantization $\Lam(\R)=\frac14\Z_{\ge0}\cup\{\infty\}$.
\end{itemize}
For a general irrational $\alpha$ whose partial quotients tend to infinity, we
First treat polynomial-growth denominator recurrences
$Q_{j+1}=(cj^r+d)Q_j+Q_{j-1}$ in
Theorems~\ref{thm:general-rec}--\ref{thm:poly-phase-transition} and then, in
Section~\ref{sec:growing-rec} (Theorems~\ref{thm:growing-rec-spectrum}
and~\ref{thm:growing-rec-dim-phase}, and
Corollary~\ref{cor:growing-rec-three-regimes}), establish the complete unified
theory for arbitrary growing-quotient recurrences $Q_{j+1}=A_jQ_j+Q_{j-1}$ with
$A_j\in\Z_{\ge2}$ and $A_j\to\infty$. In particular, Theorem~\ref{thm:growing-rec-spectrum}(i)
proves the exact asymptotic identity
$\limsup_{j\to\infty}\|Q_jx\|/(Q_j|Q_j\theta-P_j|)=\limsup_{j\to\infty}A_j\|Q_jx\|$
matching the Kraaikamp--Liardet error normalization across all growing-quotient
recurrences, while Theorem~\ref{thm:growing-rec-dim-phase} reveals how the upper
logarithmic growth exponent $\gamma_{\sup}=\limsup_{j\to\infty}\frac{\log A_j}{\log Q_j}\in[0,\infty]$
governs both the subcritical Hausdorff-dimension distortion
$(1-\tau)/(1+\tau\gamma_{\sup})$ ($0\le\tau<1$) and the supercritical secondary
collapse from $L_G$ to $\frac{1}{g}\Z$ across $\tau_{\mathrm{crit}}=1+1/\gamma_{\sup}$.
For the modular continued-fraction alphabets $F_{m,r}=\{[0;a_1,a_2,\dots]:a_n\equiv r\pmod m\}$,
the underlying infinite-alphabet dimension theory is that of Mauldin and
Urba\'nski \cite{MU96} (see \cite[Theorem~2.3]{CLU}), Hensley \cite{Hensley1996}
(whose polynomial-time paper carries a published erratum), Jenkinson and
Pollicott \cite{JP2018}, and Chousionis, Leykekhman and Urba\'nski \cite{CLU},
who proved that every arithmetic-progression alphabet has full dimension
spectrum, while Chitanga, Lemmens and Nussbaum \cite{CLN2026} recently settled
full dimension spectra $[0,\dimH F_{m,r}]$ for single arithmetic progressions
(not addressing our uniform Lipschitz invariance across $E_m$, global strict
separation $\dimH F_{m,m}<\dimH F_{m,1}$, constants $C_r$, or multi-residue
asymptotics). Asymptotic expansions of the Hausdorff dimension when generators
with small derivative are added to or removed from a conformal IFS were
developed by Das, Fishman, Simmons and Urba\'nski \cite{DFSU}, who
strengthened Hensley's expansion \cite{Hensley92} for partial quotients $\le N$
and Good's estimates \cite{Good41} for partial quotients $\ge N$. The families
$F_{m,S,B}$ fall within the scope of their perturbation theorem
\cite[Theorem~4.1]{DFSU} (fixed finite $S$, added generators $\ge m-\max B$);
in their notation the leading constant is $1/\bar c=\sum_{n\ge0}\sum_{w\in S^n}q(w)^{-1}=C_S$
when $s_0(S)<1/2$. Arithmetic-progression alphabets are not treated in
\cite{DFSU}, and Theorems~\ref{thm:residues-asymptotic}
and~\ref{thm:multi-residue-asymptotic} make the first terms explicit for them;
the $\log m/m^2$ coefficient has the same structure as Hensley's
$72/\pi^4=2(6/\pi^2)^2$ (cf.\ \cite[(1.4)]{DFSU}). Part~(i) of
Theorem~\ref{thm:multi-residue-asymptotic} is the classical Bowen formula
written as the abscissa of a Poincar\'e series; our contribution there is a
proof that avoids the pressure formalism, together with its formalization.
The upper bound in part~(v), and hence the supercritical case of part~(ii),
follow from the Heinemann--Urba\'nski estimate \cite{HU02} (in the form of
\cite[Theorem~2.4]{CLU}) combined with the Lyapunov bound
\cite[Proposition~2.6]{CLU}, and the case $S=\emptyset$ of part~(ii) follows
from Good's theorem by monotonicity. Part~(iii) is a quantitative form of the
Laurent expansion of $\zeta$ at $1$ \cite[(25.2.4)]{DLMF}. A targeted negative
search (arXiv, Crossref and web search; coverage recorded in the verification
archive) found no prior appearance of $C_r$, of $\psi_{\mathrm{Fib}}$, or of the
coefficient $\frac12 d^2C_S(\gamma C_S-C_S')$ as asymptotic dimension constants
of residue-class digit sets; this does not exclude such an appearance.
% Theorem/equation numbers cited for CLU and DFSU above are those of arXiv:1805.11904v1 and
% arXiv:2007.10554v2; the numbering of the published versions has not been verified
% (verification/paper2_novelty_review_20261125.md).
Accordingly, we
regard Theorems~\ref{thm:class} and~\ref{thm:gap} as a variant of known results.
The main new contributions are the complete spectrum discretization and
three-way positive-level splitting (Theorems~\ref{thm:full-spectrum},
\ref{thm:trichotomy}, \ref{thm:general-rec}--\ref{thm:poly-phase-transition},
and~\ref{thm:growing-rec-spectrum}--\ref{thm:growing-rec-dim-phase}), the critical gauge theorems
(Theorem~\ref{thm:gauge} and Corollary~\ref{cor:exact-gauge}), the machine-checked
dimension-zero and subcritical dimension theorems (Theorems~\ref{thm:finite},
\ref{thm:upper}, and~\ref{conj:weight}), the continued-fraction window and
residue theorems (Theorems~\ref{thm:windows}--\ref{thm:residues-general} and the
strict separation in Theorem~\ref{thm:residues-asymptotic}(ii)), and explicit
residue-class instances of known asymptotic methods
(Theorems~\ref{thm:residues-asymptotic}(iii)--(iv)
and~\ref{thm:multi-residue-asymptotic}).

Related metric questions for Cantor series expansions, where the approximation
is measured by the tail $x-\omega_n(x)$ of the partial sums, were studied by
Cheng and Zhang~\cite{CZ}, who determine Hausdorff dimensions of the
corresponding limsup sets, and recent exact-approximation frameworks by Baker
and Ward~\cite{BW2025} and Tan, Tian, Wang and Wu~\cite{TTWW2026} establish
quantitative criteria and infinite Hausdorff $f$-measures for exact
approximation sets in regular spaces; discrete spectral gaps, signed-coset
collapse, and exact recurrence gauges for non-autonomous super-exponential
denominators $(Q_j)$ are not considered there.
For the selected sequence $(Q_j)$, Corollary~8 of Bugeaud, Liao and Rams
\cite{BLR} applies directly since $Q_{j+1}/Q_j\to\infty$. Theorems~\ref{thm:poly-phase-transition} and
\ref{thm:growing-rec-dim-phase} are compared with \cite{BLR} in
Remarks~\ref{rem:poly-BLR} and~\ref{rem:growing-BLR}. Corollary~8 gives dimension
one for every constant target $c\in[0,1/2]$ at weight $\alpha=0$ in
Theorem~\ref{conj:weight} (in particular for $\{x:\|Q_jx\|\to0\}$), by taking
the centres equal to $c$ and using $|\|y\|-\|c\||\le\|y-c\|$; integer
translation invariance transfers full dimension between $\R$ and $[0,1]$. For
nonzero weights $0<\alpha<1$, the logarithmic growth condition (1.6) of their
Theorem~7 fails for the linear maps $x\mapsto Q_jx$ because
$\log Q_{j+1}/\log Q_j\to1$, and classical mass-transference methods for
$\limsup$ shrinking-target sets \cite{BV} do not apply to the eventual
(liminf-type) bounds defining $E_c$ and the prescribed-limit sets. We establish
the exact dimension $1-\alpha$ here by an elementary covering upper bound and a
shifted-grid Moran construction.

Regarding Hausdorff gauges at critical levels, Cabrelli, Mendivil, Molter and
Shonkwiler \cite[\S5, Theorem~3]{CMMS} construct a gauge giving positive
\emph{finite} measure to a Cantor set from its gap sequence, whereas Moreira,
Erazo and Angelini \cite{MEA} show that the classical Lagrange level at value
$3$ has either zero or non-$\sigma$-finite measure for every dimension function.
Here the recurrence specifies the gauge $h_k(1/Q_n)=(2k+1)^{-n}$, and the entire
critical level $\{\Lam=k/4\}$ on $\R/\Z$ has \emph{infinite, $\sigma$-finite}
$\mathcal H^{h_k}$ measure, proved by combining a countable decomposition by
initial recurrence data with translation invariance under the dense subgroup
$L$.

Relative to earlier literature, our setting has four distinguishing features:
(i) a \emph{subsequence} of convergent denominators of $\ee$ (every third),
(ii) the critical \emph{weight} $j$, (iii) the lattice
$L=\Z+\Z\frac{\ee+1}{2}$ (note $(\ee+1)/2$, not $\ee$), and (iv) the
\emph{sharp} quantum $1/4$. The weight $j$ is \emph{critical} for the covering
argument: it makes the per-parent branching number uniformly bounded, yielding
dimension $0$ for \emph{all} finite level sets (Theorem~\ref{thm:finite}),
whereas without the weight the branching grows linearly and the unweighted
zero-set $\{x:\|Q_jx\|\to0\}$ has full dimension $1$, while the weighted
zero-set $\{x:j\|Q_jx\|\to0\}$ equals the countable group $L$.

